\documentclass[12pt]{article}
\usepackage[margin=1in]{geometry}
\usepackage{setspace}
\usepackage{booktabs}
\usepackage{amsmath}
\usepackage{natbib}
\usepackage{hyperref}
\usepackage{caption}
\usepackage{array}
\usepackage{dcolumn}
\usepackage{rotating}
\usepackage{threeparttable}

\doublespacing
\bibpunct{(}{)}{;}{a}{,}{,}

\title{Women Finance Ministers' Political Affiliation and Fiscal Policy Outcomes}
\author{Jhinuk Banerjee\\
        University of Oklahoma}
\date{April 2026}

\begin{document}

\maketitle

%% -----------------------------------------------------------
\begin{abstract}
This paper examines whether women Finance Ministers' political affiliation shapes
their fiscal policy choices, and whether those effects are conditioned by the
ideological alignment or distance between the Finance Minister and the head of
government. Using a balanced country-year panel of Finance Ministers drawn from
WhoGov (within and cross-sectional versions merged), matched with ideology data
from the Global Leader Ideology Dataset (GLID) and fiscal outcomes from the IMF
World Economic Outlook, I estimate a two-way fixed effects (TWFE) model with
country and year fixed effects across up to 177 countries from 1980--2023.
The analysis identifies 411 women Finance Ministers across the full sample.
Causal identification addresses the non-random ``glass cliff'' appointment
pattern by instrumenting for women Finance Minister appointment using the count
of women in other cabinet positions. The paper contributes to the literature
on gender and economic institutions, descriptive versus substantive representation,
and the political economy of fiscal policy.
\end{abstract}

\newpage

%% -----------------------------------------------------------
\section{Introduction}
\label{sec:intro}
%% -----------------------------------------------------------

Women's representation in political institutions has grown modestly over recent
decades, but men continue to dominate economic policymaking positions globally.
As of early 2026, women hold only 22.4\% of cabinet positions worldwide, a
figure that has declined from 2024, and only 14 countries have achieved gender
parity in their cabinets. The Finance Ministry one of the most powerful and
consequential cabinet portfolios remains particularly male-dominated. Against
this backdrop, a central question for political economy and gender studies
arises: when women do break through to head Finance Ministries, do they govern
differently from men?

A large body of scholarship documents gender differences in political attitudes,
risk preferences, and distributive priorities. However, much less is known about
whether partisan identity moderates or supersedes these gender differences once
women enter executive positions of real economic authority. Is a left-wing woman
Finance Minister more similar to a left-wing man than to a right-wing woman?
Does ideological alignment between the Finance Minister and the head of government
mediate the degree to which any gender-specific fiscal patterns emerge?

This paper addresses these questions by constructing a novel country-year panel
combining three datasets: WhoGov, the Global Leader Ideology Dataset (GLID), and
the IMF World Economic Outlook (WEO). The merged dataset covers up to 177
countries from 1980--2023 and identifies 411 women Finance Ministers---a
substantially larger sample than prior cross-national work. The core empirical
strategy is a two-way fixed effects (TWFE) model augmented with an instrumental
variable (IV) to address the endogeneity of women FM appointment.

The paper makes three contributions. First, it provides the most comprehensive
cross-national analysis to date of women Finance Ministers' fiscal policy
outcomes. Second, it introduces ideology as a moderating variable, testing
whether the partisan identity of the Finance Minister shapes fiscal behaviour
differently for women than men. Third, it engages directly with the glass cliff
literature \citep{armstrong2024} by proposing and implementing an instrument
for women FM appointment to address crisis-driven selection bias.

The remainder of the paper proceeds as follows. Section~\ref{sec:lit} reviews
the relevant literature. Section~\ref{sec:data} describes the data and the
instrument. Section~\ref{sec:empirics} presents the empirical strategy.
Section~\ref{sec:results} discusses expected results.

%% -----------------------------------------------------------
\section{Literature Review}
\label{sec:lit}
%% -----------------------------------------------------------

The paper builds on three strands of literature: the relationship between
demographic identity and policymaking in economic institutions; the evidence
on whether women leaders govern differently; and the political economy of
fiscal policy and ideology.

A growing body of work establishes that the demographic composition of economic
institutions shapes both policy decisions and public perceptions of legitimacy.
\citet{mcdowell2024} show that Black representation within the Federal Reserve
increases institutional trust among Black respondents, establishing that
descriptive identity within economic institutions has direct consequences for
perceived state legitimacy. \citet{bodea2025} find that citizens on average
trust women central bankers less, particularly during high inflation, pointing
to a persistent gender bias in the perception of economic competence.
\citet{dutta2024} document structural barriers for women within the Federal
Reserve itself, finding gender gaps in publication output despite comparable
or greater work effort by women economists. \citet{charlety2016} show that
central bank board appointments follow gendered replacement dynamics---women
are significantly more likely to be appointed when replacing another woman---
suggesting a tokenism or quota-substitution logic that has direct parallels in
Finance Ministry appointments.

The evidence on whether women leaders \textit{govern differently} is mixed and
context-dependent. \citet{chattopadhyay2004} provide clean experimental evidence
from India showing that randomly assigned female village council leaders invest
differently from male leaders, particularly in public goods valued by women.
\citet{atchison2009} find that higher female cabinet representation is associated
with more female-friendly social policy across a cross-national sample. In the
fiscal and monetary domains, however, results are more nuanced.
\citet{masciandaro2020} find that greater female representation on monetary
policy boards is associated with more hawkish interest rate decisions, directly
contradicting the stereotype that women policymakers are dovish or risk-averse.
\citet{suzuki2018} document lower fiscal risk-taking associated with women in
local public finance in Japan, but note strong dependence on institutional
context. \citet{moessinger2012} establishes that personal characteristics of
Finance Ministers, including age and education, affect public debt trajectories,
confirming that minister-level heterogeneity matters for fiscal outcomes.

A critical strand for identification is the glass cliff literature.
\citet{armstrong2024} find that women are disproportionately appointed as
Finance Ministers during financial crises, and that this pattern is strongest
in countries with no prior female head of government. Crucially, crisis periods
do not reduce women's tenure relative to men, implying that the glass cliff is
a selection phenomenon rather than a performance one. \citet{kern2025} document
a parallel pattern in central banking: sovereign debt crises increase the
probability of women being appointed to central bank boards, and such
appointments are associated with more conservative monetary outcomes. Together,
this evidence implies that naive TWFE estimation of fiscal outcomes on women FM
indicators will be downward biased, since women are appointed precisely when
fiscal conditions are deteriorating.

On the ideology side, \citet{vanommeren2021} show strong and persistent
ideological effects in central banking, with right-leaning appointees pursuing
more hawkish policy throughout their tenure. This suggests that for Finance
Ministers, the partisan identity of the appointment may be as important as
gender. \citet{krook2012} and \citet{obrien2015} document the determinants of
female cabinet appointments globally, while \citet{field2021} examines the
institutional conditions under which gender shapes ministerial portfolio
allocation. None of this work, however, connects women Finance Ministers'
partisan identity to their fiscal policy choices, the gap this paper addresses.

%% -----------------------------------------------------------
\section{Data}
\label{sec:data}
%% -----------------------------------------------------------

\subsection{Dataset Construction}

The analysis draws on three primary datasets merged into a country-year panel.

\textbf{Cabinet data.} The WhoGov dataset \citep{nyrup2020} combines
within-country and cross-sectional versions to provide individual-level data on
cabinet ministers across 177 countries from 1945 to 2023. Key variables include
minister name, gender, party affiliation, portfolio assignment, and tenure dates.
Finance Ministers are identified using the binary indicator \texttt{m\_Finance},
constructed by WhoGov to flag Finance Ministers regardless of variation in
position titles across countries---an important distinction, since relying on
the position string alone would miss approximately 180 of the 411 women Finance
Ministers in the sample. The merged and balanced version of the panel contains
9,101 country-year Finance Minister observations spanning 1963--2023.

\textbf{Ideology data.} Finance Minister party ideology is constructed in two
steps. First, an original ideology variable was built for each Finance Minister
using a combination of AI-assisted coding and Wikipedia biographical sources for each political party in each country,
classifying each minister's partisan affiliation on a 1--6 left-right scale
(1 = far left, 6 = far right) across the full country-year sample. This
minister-level ideology variable (\texttt{fm\_ideology1}) is then
cross-validated against and supplemented by the Global Leader Ideology Dataset
(GLID), which provides expert-coded and survey-based left-right scores for heads
of government across countries and years. The GLID merge also yields the
head-of-government ideology variable, enabling construction of the ideological
distance between the Finance Minister and the HOG. The variable
\texttt{fm\_ideology1} is complete for all 9,101 observations, with mean 3.84
and standard deviation 1.75 on the 1--6 scale.

\textbf{Fiscal outcomes.} Fiscal data come from the IMF World Economic Outlook
(WEO) database, providing annual country-level general government revenue,
expenditure, net lending/borrowing (fiscal balance), and gross debt---all as a
share of GDP---for 196 countries from 1980--2023. The WEO is preferred over the
IMF Government Finance Statistics (GFS) because its broader coverage avoids the
selection problem that would arise from restricting to the 45--74 countries with
consistent GFS reporting, which are disproportionately high-income OECD economies
where both women Finance Ministers and constrained fiscal institutions are more
common. The WEO data were cleaned from wide to long format, with non-numeric
entries removed, yielding 8,624 usable country-year observations.

\textbf{Merge.} The WhoGov+GLID Finance Minister panel is merged many-to-one
with the WEO fiscal data on ISO3 country code and year. The merge yields 7,053
matched observations out of 9,101 total minister-year observations. The 2,048
unmatched master observations are pre-1980 years excluded from all regressions.
The effective analysis window is 1980--2023, with approximately 140--169
countries per year. After matching, 5,393 observations have non-missing fiscal
balance data (76.5\% of matched observations), forming the core analysis sample.
Summary statistics from the merged dataset are pending and will be presented in
the final version of the paper.

\subsection{Instrument: Women in Other Cabinet Positions}
\label{sec:instrument}

A central identification challenge is that women Finance Minister appointments
are not random. The glass cliff literature \citep{armstrong2024,kern2025}
establishes that women are disproportionately appointed during fiscal crises,
precisely when fiscal outcomes are already deteriorating. A naive TWFE estimate
will therefore be downward biased.

To address this, I instrument for women FM appointment using the count of women
in other cabinet positions in the same country-year:
\[
Z_{it} = \texttt{n\_female\_minister}_{it} - \mathbf{1}[\text{gender}_{it} = \text{Female}]
\]
The \textit{relevance condition} is that countries with more women in cabinet
generally have a higher probability of appointing a woman to the Finance
Ministry, through both supply-side effects (a larger pool of politically
experienced women) and norm diffusion within the executive. The
\textit{exclusion restriction} requires that the number of women in other
ministries affects fiscal outcomes only through the Finance Minister channel.
This is defensible because ministries such as Agriculture, Foreign Affairs,
and Social Affairs do not directly control the budget. The main threat to
exclusion---a progressive government simultaneously placing women everywhere
and pursuing expansionary fiscal policy---is addressed by including HOG ideology
as a control. A secondary instrument to be explored is the regional diffusion
variable: the share of countries in the same region that had a woman Finance
Minister in year $t-1$, capturing peer-pressure and demonstration effects
plausibly exogenous to any given country's fiscal balance.

%% -----------------------------------------------------------
\section{Empirical Strategy}
\label{sec:empirics}
%% -----------------------------------------------------------

\subsection{IV Two-Way Fixed Effects}
\label{sec:ivtwfe}

The baseline estimating equation is:
\vspace{-0.4cm}
\begin{equation}
\label{eq:twfe}
Y_{it} = \beta_1 \text{FemaleFM}_{it} + \beta_2 \text{Ideology}_{it}
       + \beta_3 (\text{FemaleFM}_{it} \times \text{Ideology}_{it})
       + \gamma X_{it} + \alpha_i + \lambda_t + \varepsilon_{it}
\end{equation}

where $Y_{it}$ is a fiscal outcome in country $i$ and year $t$ (revenue,
expenditure, fiscal balance, or gross debt as \% of GDP); $\text{FemaleFM}_{it}$
is an indicator equal to 1 when a woman holds the Finance Ministry;
$\text{Ideology}_{it}$ is the Finance Minister's ideology score (1--6 scale);
$X_{it}$ is a vector of time-varying controls; $\alpha_i$ are country fixed
effects; and $\lambda_t$ are year fixed effects. Standard errors are clustered
at the country level. The coefficient of interest is $\beta_3$, capturing
whether ideology shapes fiscal behaviour differently by gender.

Time-varying controls include GDP growth rate, inflation rate, a financial
crisis indicator, HOG ideology, HOG leader experience (years in office), the
democracy indicator (Regimes of the World), and the share of women in the full
cabinet. Given the glass cliff endogeneity, $\text{FemaleFM}_{it}$ is treated
as endogenous and instrumented by $Z_{it}$ (Section~\ref{sec:instrument}).
The IV-TWFE first stage is:

\begin{equation}
\label{eq:firststage}
\text{FemaleFM}_{it} = \pi_1 Z_{it} + \pi_2 \text{Ideology}_{it}
                      + \gamma X_{it} + \alpha_i + \lambda_t + \nu_{it}
\end{equation}

with instrument strength evaluated by the first-stage $F$-statistic
(threshold: $F > 10$). The second stage substitutes
$\widehat{\text{FemaleFM}}_{it}$ into equation (\ref{eq:twfe}).

\subsection{HOG Alignment Specification}
\label{sec:hogalign}

To assess whether ideological alignment between the Finance Minister and the
head of government conditions the gender-ideology effect, I extend
equation~(\ref{eq:twfe}) to:

\begin{align}
\label{eq:extended}
Y_{it} &= \beta_1 \text{FemaleFM}_{it} + \beta_2 \text{Ideology}_{it}
        + \beta_3 (\text{FemaleFM}_{it} \times \text{Ideology}_{it}) \nonumber \\
       &\quad + \beta_4 (\text{FemaleFM}_{it} \times \text{HOGIdeology}_{it})
        + \gamma X_{it} + \alpha_i + \lambda_t + \varepsilon_{it}
\end{align}

$\beta_4$ captures whether the head of government's ideology moderates the
fiscal behaviour of women Finance Ministers---for instance, whether ideological
co-partisanship between the FM and HOG grants the FM greater fiscal discretion,
or whether a misaligned HOG constrains it. An alternative specification
replaces the HOG ideology level with the ideological distance measure
$|\text{Ideology}_{it} - \text{HOGIdeology}_{it}|$ to test whether the
\textit{degree} of misalignment, not just its direction, shapes outcomes.

%% -----------------------------------------------------------
\section{Expected Results}
\label{sec:results}
%% -----------------------------------------------------------

Results are pending completion of the Stata analysis on the merged dataset.
The following expectations are grounded in the theory and empirical literature
reviewed above.

On the OLS TWFE baseline, I expect a negative and significant $\beta_1$:
women Finance Ministers associated with worse fiscal balances on average,
reflecting the glass cliff selection mechanism rather than a causal gender
effect. The ideology interaction $\beta_3$ is the central test: if ideology
dominates gender, the interaction should be significant and the pure gender
coefficient should attenuate substantially once ideology is controlled. A
positive $\beta_3$ would imply that left-wing women Finance Ministers are
associated with more expansionary fiscal policy relative to right-wing women,
consistent with standard partisan fiscal theory.

On the IV-TWFE, I expect the women FM coefficient to move toward zero or
turn positive once the glass cliff selection is purged---mirroring the
pattern in \citet{armstrong2024} where crisis-driven selection accounts for
much of the apparent gender-fiscal correlation. A strong first stage
($F > 10$) is expected given the tight mechanical relationship between
women's representation in the broader cabinet and Finance Ministry
appointments.

On the HOG alignment specification, I expect $\beta_4$ to be positive and
significant---women Finance Ministers with ideologically co-partisan HOGs
should exhibit stronger partisan fiscal signals---consistent with the
principal-agent logic that ministerial fiscal discretion is greatest when the
FM and HOG share partisan interests \citep{vanommeren2021}. These results,
taken together, would establish that the \textit{type} of woman appointed---
her ideology and her relationship to the head of government---matters more
than gender per se for fiscal outcomes.


%% -----------------------------------------------------------
\begin{thebibliography}{}
 
\bibitem[Armstrong et~al., 2024]{armstrong2024}
Armstrong, B., Barnes, T.~D., Chiba, D., and O'Brien, D.~Z. (2024).
\newblock Financial crisis and the selection and survival of women finance
  ministers.
\newblock {\em American Political Science Review}, 118(1):1--16.
 
\bibitem[Atchison and Down, 2009]{atchison2009}
Atchison, A. and Down, I. (2009).
\newblock Women cabinet ministers and female-friendly social policy.
\newblock {\em Poverty and Public Policy}, 1(2):1--23.
 
\bibitem[Bodea and Kerner, 2025]{bodea2025}
Bodea, C. and Kerner, A. (2025).
\newblock Expectations, gender bias and {Federal Reserve} talk: Do {Americans}
  trust women as central bankers?
\newblock {\em Journal of Politics}.
 
\bibitem[Charlety et~al., 2016]{charlety2016}
Charlety, P., Romelli, D., and Santacreu-Vasut, E. (2016).
\newblock Appointments to central bank boards: Does gender matter?
\newblock {\em ESSEC Working Paper}.
 
\bibitem[Chattopadhyay and Duflo, 2004]{chattopadhyay2004}
Chattopadhyay, R. and Duflo, E. (2004).
\newblock Women as policy makers: Evidence from a randomized policy experiment
  in {India}.
\newblock {\em Econometrica}, 72(5):1409--1443.
 
\bibitem[Dutta and Tzur-Ilan, 2024]{dutta2024}
Dutta, D.~D. and Tzur-Ilan, N. (2024).
\newblock Gender gaps in the {Federal Reserve} system.
\newblock {\em Finance and Economics Discussion Series, Board of Governors of
  the Federal Reserve System}.
 
\bibitem[Field, 2021]{field2021}
Field, B.~N. (2021).
\newblock Ministers, gender and political appointments.
\newblock {\em Government and Opposition}, 56(2):229--247.
 
\bibitem[Kern et~al., 2025]{kern2025}
Kern, A., Reinsberg, B., and Romelli, D. (2025).
\newblock Empowering women in central banking.
\newblock {\em Journal of European Public Policy}, 32(1):1--25.
 
\bibitem[Krook and O'Brien, 2012]{krook2012}
Krook, M.~L. and O'Brien, D.~Z. (2012).
\newblock All the president's men? {The} appointment of female cabinet
  ministers worldwide.
\newblock {\em Journal of Politics}, 74(3):840--855.
 
\bibitem[Masciandaro et~al., 2020]{masciandaro2020}
Masciandaro, D., Profeta, P., and Romelli, D. (2020).
\newblock Do women matter in monetary policy boards?
\newblock {\em Working Paper}.
 
\bibitem[McDowell and Steinberg, 2024]{mcdowell2024}
McDowell, D. and Steinberg, D.~A. (2024).
\newblock Black representation and the popular legitimacy of the {Federal
  Reserve}.
\newblock {\em European Journal of Political Economy}.
 
\bibitem[Moessinger, 2012]{moessinger2012}
Moessinger, M.-D. (2012).
\newblock Do personal characteristics of finance ministers affect the
  development of public debt?
\newblock {\em ZEW Discussion Paper No. 12-068}.
 
\bibitem[Nyrup and Bramwell, 2020]{nyrup2020}
Nyrup, J. and Bramwell, S. (2020).
\newblock Who governs? {A} new global dataset on members of cabinets.
\newblock {\em American Political Science Review}, 114(4):1366--1374.
 
\bibitem[O'Brien et~al., 2015]{obrien2015}
O'Brien, D.~Z., Mendez, M., Peterson, J.~C., and Shin, J. (2015).
\newblock Letting down the ladder or shutting the door: Female prime ministers,
  party leaders, and cabinet ministers.
\newblock {\em Politics and Gender}, 11(4):689--717.
 
\bibitem[Suzuki and Avellaneda, 2018]{suzuki2018}
Suzuki, K. and Avellaneda, C.~N. (2018).
\newblock Women and risk-taking behaviour in local public finance.
\newblock {\em Public Management Review}, 20(12):1741--1767.
 
\bibitem[van Ommeren and Piccillo, 2021]{vanommeren2021}
van Ommeren, E. and Piccillo, G. (2021).
\newblock The central bank governor and interest rate setting by committee.
\newblock {\em CESifo Economic Studies}, 67(2):134--167.
 
\end{thebibliography}
%% -----------------------------------------------------------
 
\end{document}